% Einheit 3 von 4 – Flächen- und Volumeneinheiten (bank/einheiten/e3.jsonl, zusammenbau v0.1)
%% TODO Zweigzeile Teil 1: was hier gelernt wird (Satz in Schülersprache aus dem Einheitstitel) – Einheit 3
\einheitenkopf[e3]{Einheit 3 von 4 · Flächen- und Volumeneinheiten}
\zweigzeile{ab Kl. 5 · P10}
\setcounter{aufgabe}{28}

%% TODO \verfahren{…}: Verfahrensüberschrift in Sachsprache fehlt in der Bank (2.3 b)
%% TODO Titel: Ich-kann-Satz fehlt in der Bank (Platzhalter: Kettenname) – Nr. 29
%% TODO Anweisung: ein Satz über der ersten Teilaufgabe fehlt in der Bank – Nr. 29
\begin{aufgabe}{Flächen- und Volumeneinheiten umrechnen}
\begin{teile}
\teil Quadratmeter und Quadratdezimeter – welche Umrechnungszahl? Kreuze an.\\ \kreuz{zehn}\\ \kreuz{hundert}\\ \kreuz{tausend}
\end{teile}
%% TODO Darstellung für schwach fehlt in der Bank (einheiten-e3-k1-s1-v1; Abschnitt „Für schwache Schüler“)
\swz{$3$ m² – wie viel dm²?}{}
%% TODO Darstellung für schwach fehlt in der Bank (einheiten-e3-k1-s1-v2; Abschnitt „Für schwache Schüler“)
\swz{$7$ cm² – wie viel mm²?}{}
%% TODO Darstellung für schwach fehlt in der Bank (einheiten-e3-k1-s1-v3; Abschnitt „Für schwache Schüler“)
\swz{$500$ dm² – wie viel m²?}{}
%% TODO Darstellung für schwach fehlt in der Bank (einheiten-e3-k1-s1-v4; Abschnitt „Für schwache Schüler“)
\swz{$12$ dm² – wie viel cm²?}{}
%% TODO Darstellung für schwach fehlt in der Bank (einheiten-e3-k1-s1-v5; Abschnitt „Für schwache Schüler“)
\swz{$900$ mm² – wie viel cm²?}{}
\swfrage{Was bleibt gleich, was ändert sich?}
%% TODO Darstellung für schwach fehlt in der Bank (einheiten-e3-k1-s2-v1; Abschnitt „Für schwache Schüler“)
\swz{$4$ m² – wie viel cm²?}{}
%% TODO Darstellung für schwach fehlt in der Bank (einheiten-e3-k1-s3-v1; Abschnitt „Für schwache Schüler“)
\swz{$7$ a – wie viel m²?}{}
%% TODO Darstellung für schwach fehlt in der Bank (einheiten-e3-k1-s4-v1; Abschnitt „Für schwache Schüler“)
\swz{Ordne, mit der kleinsten beginnend: m³, mm³, dm³, cm³}{}
%% TODO Darstellung für schwach fehlt in der Bank (einheiten-e3-k1-s5-v1; Abschnitt „Für schwache Schüler“)
\swz{$5$ dm³ – wie viel cm³?}{}
%% TODO Darstellung für schwach fehlt in der Bank (einheiten-e3-k1-s6-v1; Abschnitt „Für schwache Schüler“)
\swz{$6$ l – wie viel dm³?}{}
%% TODO Darstellung für schwach fehlt in der Bank (einheiten-e3-k1-s7-v1; Abschnitt „Für schwache Schüler“)
\swz{$1{,}8$ l – wie viel ml?}{}
\begin{teile}
\teil Was bedeutet „Kilo“ in Kilowatt? Kreuze an.\\ \kreuz{tausend}\\ \kreuz{hundert}\\ \kreuz{ein Tausendstel}
\end{teile}
%% TODO Darstellung für schwach fehlt in der Bank (einheiten-e3-k1-s9-v1; Abschnitt „Für schwache Schüler“)
\swz{Eine Flasche enthält $1{,}25$ l Saft. Wie viele Gläser zu $200$ ml kann man damit ganz füllen? \hfill (P10 2022 OS)}{}
\swz[3]{Eine Pumpe fördert $12\,500$ l pro Stunde. Wie lange braucht sie, um einen Teich mit $75$ m³ Wasser zu füllen? \hfill (P10 2017 OS)}{}
\end{aufgabe}

%% TODO \verfahren{…}: Verfahrensüberschrift in Sachsprache fehlt in der Bank (2.3 b)
%% TODO Titel: Ich-kann-Satz fehlt in der Bank (Platzhalter: Kettenname) – Nr. 30
%% TODO Anweisung: ein Satz über der ersten Teilaufgabe fehlt in der Bank – Nr. 30
\begin{aufgabe}{Flächen- und Volumeneinheiten im Koordinatensystem, Sek II}
\begin{teile}
\teil Eine Längeneinheit entspricht $15$ cm; gesucht ist der Flächeninhalt eines Schildes. Was wird umgerechnet? Kreuze an.\\ \kreuz{Länge}\\ \kreuz{Fläche}\\ \kreuz{Volumen}
\end{teile}
%% TODO Darstellung für schwach fehlt in der Bank (einheiten-e3-k2-s1-v1; Abschnitt „Für schwache Schüler“)
\swz{Eine Fläche ist $6$ Flächeneinheiten groß; eine Längeneinheit entspricht $5$ cm. Wie viel cm²?}{}
%% TODO Darstellung für schwach fehlt in der Bank (einheiten-e3-k2-s1-v2; Abschnitt „Für schwache Schüler“)
\swz{Eine Fläche ist $3$ Flächeneinheiten groß; eine Längeneinheit entspricht $30$ cm. Wie viel cm²?}{}
%% TODO Darstellung für schwach fehlt in der Bank (einheiten-e3-k2-s1-v3; Abschnitt „Für schwache Schüler“)
\swz{Eine Fläche ist $12$ Flächeneinheiten groß; eine Längeneinheit entspricht $2$ cm. Wie viel cm²?}{}
%% TODO Darstellung für schwach fehlt in der Bank (einheiten-e3-k2-s1-v4; Abschnitt „Für schwache Schüler“)
\swz{Eine Fläche ist $7{,}5$ Flächeneinheiten groß; eine Längeneinheit entspricht $4$ cm. Wie viel cm²?}{}
%% TODO Darstellung für schwach fehlt in der Bank (einheiten-e3-k2-s1-v5; Abschnitt „Für schwache Schüler“)
\swz{Eine Fläche ist $9$ Flächeneinheiten groß; eine Längeneinheit entspricht $15$ cm. Wie viel cm²?}{}
\swfrage{Was bleibt gleich, was ändert sich?}
%% TODO Darstellung für schwach fehlt in der Bank (einheiten-e3-k2-s2-v1; Abschnitt „Für schwache Schüler“)
\swz[3]{Die Vorderseite einer Gartenhütte ist ein Rechteck, $3$ Längeneinheiten breit und $2$ hoch, mit einem Dachdreieck darauf (Grundseite $3$, Höhe $1{,}5$). Eine Längeneinheit entspricht einem Meter, die Hütte ist $4$ m tief. Querschnitt und Volumen?}{}
%% TODO Darstellung für schwach fehlt in der Bank (einheiten-e3-k2-s3-v1; Abschnitt „Für schwache Schüler“)
\swz[3]{Ein Aufsteller hat die Form der Fläche zwischen $g(x) = -0{,}25x^2 + 2$ (oben) und $f(x) = 0{,}25x^4 - x^2 + 1$ (unten); die Schnittstellen liegen bei $x = -2$ und $x = 2$, eine Längeneinheit entspricht $30$ cm. Weise nach, dass die Fläche $4{,}8$ Flächeneinheiten groß ist. Der Aufsteller wird auf beiden Seiten beklebt. Wie viel Folie braucht man in cm² und in m²? \hfill (FHR 2022)}{}
\end{aufgabe}

%% TODO Titel: Ich-kann-Satz fehlt in der Bank (Platzhalter: Kettenname) – Nr. 31
%% TODO Anweisung: ein Satz über der ersten Teilaufgabe fehlt in der Bank – Nr. 31
\begin{aufgabe}{Flächeneinheiten ordnen (mm² bis km², a und ha einordnen)}
%% TODO Darstellung für schwach fehlt in der Bank (einheiten-e3-k3-s1-v1; Abschnitt „Für schwache Schüler“)
\swz{Ordne, mit der kleinsten beginnend: m², ha, cm², a, km²}{}
\end{aufgabe}

%% TODO Titel: Ich-kann-Satz fehlt in der Bank (Platzhalter: Kettenname) – Nr. 32
%% TODO Anweisung: ein Satz über der ersten Teilaufgabe fehlt in der Bank – Nr. 32
\begin{aufgabe}{Repräsentanten zuordnen}
%% TODO Darstellung für schwach fehlt in der Bank (einheiten-e3-k4-s1-v1; Abschnitt „Für schwache Schüler“)
\swz{Ordne zu: Briefmarke, Kinderzimmer, Fußballfeld – $7\,000$ m²; $6$ cm²; $18$ m²}{}
\end{aufgabe}

%% TODO Titel: Ich-kann-Satz fehlt in der Bank (Platzhalter: Kettenname) – Nr. 33
%% TODO Anweisung: ein Satz über der ersten Teilaufgabe fehlt in der Bank – Nr. 33
\begin{aufgabe}{Flächen- und Volumeneinheiten umrechnen – Fehler finden, Begründen, Anwendung, Darstellungswechsel}
\swz[3]{Nils schreibt: $7$ dm² $= 70$ cm². Wo ist der Fehler?}{}
\swz[3]{Eine Längeneinheit entspricht $25$ cm, eine Fläche ist $8$ Flächeneinheiten groß. Ole rechnet: $8 \cdot 25 = 200$ cm². Wo ist der Fehler?}{}
\swfrage{Begründe: Warum passen in einen Quadratmeter hundert Quadratdezimeter und nicht zehn?}
\swfrage{Eine Längeneinheit entspricht $3$ cm. Begründe, warum eine Flächeneinheit nicht $3$ cm², sondern $9$ cm² entspricht.}
\swz[3]{Ein Aquarium ist innen $90$ cm lang, $40$ cm breit und $45$ cm hoch. Es wird bis $5$ cm unter den Rand gefüllt. Wie viele Liter Wasser braucht man?}{}
\swa{Trage $3{,}0715$ m² in die Einheitentafel ein (zwei Stellen je Einheit) und gib die Fläche in cm² an. \leerfeld[cm²]}{\sachtabelle{ccc}{m² & dm² & cm²}{\leerzelle & \leerzelle & \leerzelle}}
\end{aufgabe}

\uebersichtskasten{Flächen: Eine Fläche ist Länge mal Länge – darum ist die Umrechnungszahl 100. 1 m² = 100 dm² = 10 000 cm². Dazu 1 a = 100 m², 1 ha = 100 a. \\ Volumen: Ein Volumen ist Länge mal Länge mal Länge – darum ist die Umrechnungszahl 1000. 1 m³ = 1000 dm³; 1 dm³ = 1000 cm³. \\ Hohlmaße: 1 l = 1 dm³ und 1 ml = 1 cm³, also 1 m³ = 1000 l. \\ Vorsätze: Milli heißt ein Tausendstel, Zenti ein Hundertstel, Kilo tausend (Schreibweise mit Zehnerpotenzen → potenzen-wurzeln.md, Einheit 2). \\ Maßstab bei Flächen und Volumen (Sek II): Entspricht eine Längeneinheit 20 cm, so entspricht eine Flächeneinheit 20 cm · 20 cm = 400 cm² und eine Volumeneinheit 20 cm · 20 cm · 20 cm = 8000 cm³ – der Faktor wird mitquadriert beziehungsweise hoch drei genommen, dann erst in Quadrat- oder Kubikmeter.}
